\chapter*{Abstract}
\addcontentsline{toc}{chapter}{Abstract}

A natural language autoencoder (NLA) explains an activation inside a language model by writing a paragraph
about it, and is trained so that a second model can rebuild the activation from the paragraph. The
explanations read well, but they sometimes name the wrong specific: a city, person or instrument that was
never in the text. This report measures that failure on the released NLA for Qwen2.5-7B, using 6{,}000
Wikipedia passages about 60 entities, and asks how often it happens, why, and what removes it.

Four findings come out. First, specifics fade from the explanation faster than from the activation: sixteen
tokens after an entity, the explanation names it 7.5\% of the time while a linear probe still reads it 35\%
of the time, and 22\% of explanations contain a wrong name of the same type. Second, the training objective
barely notices: editing a wrong name out changes the reconstruction cosine by about 0.002. Third, agreement
between sampled explanations separates true names from wrong ones; keeping only names that two of four
samples share removes 74\% of wrong names and keeps 77\% of true ones, with no training, and most of that
signal is already present in one explanation's token probabilities. Fourth, the explanation survives
intervention: fine-tuning on reconstruction alone changes nothing, a cost on unsupported names cuts wrong
names by 8\% on held-out passages, and rewarding the true name does not bring it back.

The practical advice is simple: trust the gist of an NLA explanation, and do not trust a name unless
several samples agree on it.
